\section{Conclusion}
\label{sec:conclusion}

In this paper, we studied 3D scene reconstruction from noisy multi-view inputs. Unlike classical NeRF and 3D Gaussian Splatting-based methods that assume clean input images, we focused on a setting that more closely reflects real-world acquisition conditions, where multiple noise types, varying noise intensities, and cross-view consistency all play important roles.

To address this problem, we proposed \emph{DenoiseSplat}, a feed-forward 3D Gaussian splatting network tailored to noisy inputs and built upon the MVSplat framework. We constructed a multi-noise, scene-consistent noisy--clean paired multi-view dataset on RealEstate10K by injecting Gaussian, Poisson, speckle, and salt-and-pepper noise into the 2D RGB domain, and trained DenoiseSplat end-to-end to map noisy multi-view images to clean renderings without accessing any ground-truth 3D supervision. Within this framework, a dual-branch Gaussian head decouples geometry- and appearance-related parameters, enabling the network to perform geometry reconstruction and noise suppression \emph{inside} the 3D representation. At test time, a single forward pass from noisy multi-view inputs yields a clean 3DGS scene and high-quality reconstructed views, without any test-time optimization.

On the constructed noisy RE10K benchmark, we compared DenoiseSplat against the original MVSplat and a strong two-stage baseline that combines a state-of-the-art 2D denoiser with MVSplat. Experimental results show that DenoiseSplat achieves consistently better or competitive PSNR, SSIM, and LPIPS across diverse noise types and intensities, while offering better texture preservation and geometric consistency in novel view synthesis. Ablation studies further confirm that jointly training under multiple noise types and varying noise intensities, together with the proposed dual-branch Gaussian head, is crucial for obtaining robust 3DGS representations under noisy inputs.

Despite these promising results, our work still has several limitations that merit further exploration. First, our current noise modeling is based mainly on synthetic noise and does not yet capture more complex degradations such as real camera noise, motion blur, and compression artifacts. Second, our experiments are primarily conducted on RE10K, and the cross-dataset and real-world generalization of the model remains to be systematically evaluated. Future work could incorporate real-noise captures or more accurate camera noise models, extend the framework to broader degradation types and dynamic scenes, and integrate our approach with higher-level semantic cues, aiming for feed-forward 3D representations that are not only noise-robust but also better aligned with downstream 3D understanding tasks.



